\section{Appendix D. Interval and observational proofs}

The interval and observational comparisons use two features of the constructive analysis: a deterministic bound on squared-error risk and a parameter-independent transformation preserving the target. The decision-theoretic results in \cref{obj:lem:randomized-kernel-barycenter,obj:lem:randomized-kernel-tv-comparison} supply the corresponding comparisons for procedures using auxiliary randomization.

\subsection{Coverage and expected interval length}

Fix \(\alpha\in(0,1)\), and suppose that \(n,d\ge1\), \(0<q\le1\), and \cref{obj:ass:rare-arrival-slice} holds. The deterministic envelope \leanref{sym:\mathfrak U_{n,d,q}}{\ensuremath{\mathfrak U_{n,d,q}}} in \cref{obj:thm:unrestricted-mixed-count-upper} controls squared-error risk uniformly over the model in \cref{obj:def:model-class}. Its role in \cref{obj:def:risk-envelope-interval} is to determine the radius before observing the sample. Markov's inequality converts this risk control into coverage at least \(1-\alpha\). Because the target belongs to \([-1,1]\), intersecting the interval with that range preserves coverage.

The mixed-count estimator \leanref{sym:\widehat\tau^{\mathrm{mix}}_{n,d,q}}{\ensuremath{\widehat\tau^{\mathrm{mix}}_{n,d,q}}} is defined in \cref{obj:def:mixed-count-estimator}. The resulting interval \leanref{sym:I^{\mathrm{mix}}_\alpha}{\ensuremath{I^{\mathrm{mix}}_\alpha}} has length bounded by
\[
\operatorname{length}(I^{\mathrm{mix}}_\alpha)
\le
\min\left\{2,\,
2\sqrt{\frac{\mathfrak U_{n,d,q}}{\alpha}}\right\}.
\]
Thus the envelope comparison with the rate \leanref{sym:r_{n,d,q}}{\ensuremath{r_{n,d,q}}} in \cref{obj:def:frontier-rate} supplies the upper expected-length comparison in \cref{obj:thm:connected-interval-frontier}. The upper constant \leanref{sym:\overline C_\alpha}{\ensuremath{\overline C_\alpha}} may depend on the fixed miscoverage level and is uniform over the sample sizes, dimensions, and arrival floors satisfying these conditions.

The lower comparison concerns every endpoint procedure admitted by \cref{obj:def:interval-length-risk}. Its sampling component is supplied directly by \cref{obj:lem:one-cell-testing-family}. For each fixed \(\alpha\), that result supplies a finite grid of distinct target values within the same one-category family and bounds maximal expected length below at order
\[
\min\{1,(nq)^{-1/2}\}.
\]
The underlying family is common across coverage levels; the grid and its comparison constant may depend on \(\alpha\). This finite-family formulation supports every coverage level \(1-\alpha\in(0,1)\).

The many-category component of \cref{obj:thm:connected-interval-frontier} uses the separation and observation-law comparison associated with \cref{obj:def:activation-lower-handle,obj:lem:moment-matched-reciprocal-priors}. Moment matching controls statistical distinguishability, while reciprocal separation determines the target scale. Connectedness supplies the geometric link to expected length: an interval containing separated target values spans their separation. Finite testing families at the specified coverage level retain this link throughout the coverage range in the theorem.

Auxiliary endpoint randomization enters through the joint observation--action laws defined in \cref{obj:lem:randomized-kernel-tv-comparison}. Apply that result with the action space equal to the ordered endpoint pairs and the common decision rule equal to the interval procedure. The joint-law total variation equals the total variation of the observation laws. Consequently, coverage events can be compared while retaining both the observed sample and the reported endpoints. Expected length is evaluated under this same joint law, exactly as required by \cref{obj:def:interval-length-risk}.

Together, these components give the lower constant \leanref{sym:\underline c_\alpha}{\ensuremath{\underline c_\alpha}} in \cref{obj:thm:connected-interval-frontier}. The minimax expected length \leanref{sym:\mathfrak L_{n,d,q,\alpha}}{\ensuremath{\mathfrak L_{n,d,q,\alpha}}} is therefore comparable to \(\sqrt{r_{n,d,q}}\), with constants depending on \(\alpha\). The square-root scale reflects both sources of uncertainty: effective sample size and the number of covariate categories. Along sequences satisfying the rare-arrival restriction, at fixed \(\alpha\), optimal expected length tends to zero exactly when \(r_{n,d,q}\) tends to zero.

\subsection{Synthetic assignment and observational risk}

Fix a strict positivity margin \(0<\epsilon<1/2\). The observational class \leanref{sym:\mathcal D_{d,\epsilon}}{\ensuremath{\mathcal D_{d,\epsilon}}} is defined in \cref{obj:def:zeng-discrete-ate-class}, and its target is the contrast in \cref{obj:def:zeng-ate-functional}. The transformation in \cref{obj:lem:synthetic-randomization-kernel}, with arrival floor \(q=\epsilon\), assigns an independent fair arm to each observational record and retains the outcome when that arm agrees with the observational treatment. Positivity supplies the arrival floor in both arms.

The induced full-data law \(P_Z^*\), defined by \cref{obj:lem:synthetic-randomization-kernel}, belongs to the unrestricted randomized class and preserves the target:
\[
\tau(P_Z^*)=\theta_Z(P_Z).
\]
The same statement identifies the synthetic sample law with the product observed-record law of this completion. These two properties make the unrestricted risk certificate in \cref{obj:thm:unrestricted-mixed-count-upper} applicable to the mixed-count estimator evaluated on the synthetic records.

The observational estimator \leanref{sym:\widehat\theta^{\mathrm{poly}}_{n,d,\epsilon}}{\ensuremath{\widehat\theta^{\mathrm{poly}}_{n,d,\epsilon}}} in \cref{obj:thm:zeng-fixed-positivity-polynomial-frontier} averages over the independent synthetic assignments conditional on the observational sample. The barycenter comparison in \cref{obj:lem:randomized-kernel-barycenter} weakly reduces squared-error risk under this averaging. Accordingly, the constructive bound has a universal upper constant and retains the positivity dependence explicitly through
\[
r_{n,d,\epsilon}
=
\min\left\{
1,\frac{1}{n\epsilon}
+
\left[
\frac{d}{n\epsilon\log(e+n\epsilon)}
\right]^2
\right\}.
\]
Synthetic assignment converts observational treatment imbalance into selective outcome arrival; conditional averaging then removes the variation introduced by that assignment. Target preservation connects both operations to the same observational contrast.

The matching lower comparison invokes \citet[Theorem~2, Section~3.2, p.~9, together with its proof in Section~C.5, pp.~39--44]{ZengBalakrishnanHanKennedy2026}. The relevant published conclusion concerns fixed strict positivity, sufficiently large \(n\), and \(1\le d\le c_0n\log n\), where \(c_0>0\) is a constant in that conclusion. Its proof uses the submodel with zero control regression on occupied categories and supplies the lower order
\[
\frac1n+\left(\frac{d}{n\log n}\right)^2.
\]
This conclusion supplies the observational lower-bound input. The matched comparison in \cref{obj:thm:zeng-fixed-positivity-polynomial-frontier} covers every \(n,d\ge1\), incorporates bounded-risk truncation, and uses \(\log(e+n)\).

For each fixed \(\epsilon\), the observational minimax risk \leanref{sym:\mathfrak R^Z_{n,d,\epsilon}}{\ensuremath{\mathfrak R^Z_{n,d,\epsilon}}} consequently has the two-sided comparison stated in \cref{obj:thm:zeng-fixed-positivity-polynomial-frontier}, with positive finite constants \(c_\epsilon\) and \(C_\epsilon\) depending on that margin. This dependence is distinct from the universal constant in the constructive bound expressed through \(r_{n,d,\epsilon}\).

\subsection{Fixed arrival and category-growth consequences}

The same transformation determines the direction of the fixed-arrival experiment comparison. Fix \(q\in(0,1/2)\). For every observational law in \(\mathcal D_{d,q}\), \cref{obj:lem:synthetic-randomization-kernel} supplies a target-preserving completion in the unrestricted class \leanref{sym:\mathcal M^{+}_{n,d,q}}{\ensuremath{\mathcal M^{+}_{n,d,q}}} of \cref{obj:def:unrestricted-arrival-model-class}. Composing any randomized estimation procedure with the synthetic assignment gives an observational procedure with the same risk at the corresponding completion. Hence the unrestricted randomized minimax risk \leanref{sym:\mathfrak R^{+}_{n,d,q}}{\ensuremath{\mathfrak R^{+}_{n,d,q}}} dominates the observational minimax risk at positivity margin \(q\). The published observational lower comparison therefore transfers to this randomized class.

Combined with the constructive upper comparison, this transfer supports \cref{obj:thm:unrestricted-fixed-q-minimax-frontier}. Its constants \(\underline c(q)\) and \(\overline C(q)\) depend only on the fixed arrival floor and apply uniformly over \(n,d\ge1\). The resulting category-growth conclusions are recorded in \cref{obj:prop:unrestricted-fixed-q-phase-boundaries}: consistency is equivalent to
\[
d_n=o\{n\log(e+n)\},
\]
and parametric squared-error order is equivalent to
\[
d_n=O\{\sqrt n\log(e+n)\}.
\]
For an eventual bound \(d_n\le M\sqrt n\log(e+n)\), the lower risk constant may depend on \(q\), while the upper constant may depend on both \(q\) and \(M\). The eventual sample-size threshold may depend on the category sequence. These dependencies specify the fixed-arrival scope of the matched comparison and its sequence-level consequences.
